% ===== 原文 pp.18--22 =====
% [原文 p.18]
\begin{figure}[htbp]
\centering
\includegraphics[width=0.95\textwidth]{fig5.png}
\figtitle{5}{把一个 $p$-形式对称性的拓扑算符 $U_g$ 移过 $p$ 维算符 $\cO(M_p)$，会在 $\cO(M_p)$ 上留下一个拓扑局域算符 $\cO(x)$，其取值为一个数 $\phi(g)\in\C^\times$。}
\end{figure}

\begin{notebox}{关于高阶荷的注记}
这里我们遵循 [4]，只讨论了 $p$-形式对称性对维数为 $q=p$ 的算符的作用。更一般地，$p$-形式对称性也作用于维数 $q>p$ 的算符。这类作用最近在文献 [54,55] 中得到研究。遵循该处引入的术语，我们说 $q$ 维算符在 $p$-形式对称群 $G^{(p)}$ 下按 \textbf{q-荷}（q-charges）变换。

这对 0-形式对称性同样成立，此时我们有如下陈述
\begin{equation}
  \text{$G^{(0)}$ 的 $q$-荷} = \text{群 $G^{(0)}$ 的 $(q+1)$-表示}\,, \tag{2.61}
\end{equation}
其中 $(q+1)$-表示是群 $G^{(0)}$ 的熟悉的通常表示的范畴式推广。按这一术语，1-表示就是 $G^{(0)}$ 的通常表示，于是我们恢复众所周知的事实：局域算符按 0-形式对称群 $G^{(0)}$ 的通常表示变换。高阶表示（$q\ge1$）的描述超出本讲义的范围，但我们把感兴趣的读者引向文献 [54--59]，以获取物理严格程度的高阶表示导引。

类似地，对 $p\ge1$-形式对称性，我们有如下陈述
\begin{equation}
  \text{$G^{(p)}$ 的 $q$-荷} = \text{$(p+1)$-群 $\mathbb{G}^{(p+1)}_{G^{(p)}}$ 的 $(q+1)$-表示}\,, \tag{2.62}
\end{equation}
其中上述 $(p+1)$-群由群 $G^{(p)}$ 唯一确定。请注意，高阶群是通常群的范畴式推广。更多细节请读者参阅 [54,55]（另见 [60,61]，其中把高阶荷推广到了不可逆对称性）。
\end{notebox}

% [Tutor Note: p.18 (2.62) 中的 $\mathbb{G}^{(p+1)}_{G^{(p)}}$ 是双线体（blackboard bold）G，与 (2.42) 的 $G^{(p)}$ 不是同一层级的对象：$G^{(p)}$ 是普通的 $p$-形式对称群，而 $\mathbb{G}^{(p+1)}_{G^{(p)}}$ 是由它唯一确定的 $(p+1)$-群（高阶群）。核对办法：$p=0$ 时 (2.62) 应退化为 (2.61)，此时 $(p+1)$-群即 $\mathbb{G}^{(1)}_{G^{(0)}}$，其 1-表示恰是 $G^{(0)}$ 的通常表示。]

% [原文 p.19]
\begin{figure}[htbp]
\centering
\includegraphics[width=0.62\textwidth]{fig6.png}
\figtitle{6}{$p\ge1$-形式对称性的环绕（linking）作用。插入在 $S^{d-p-1}$ 上的对称性生成元 $U_g$ 作用在插入于 $M_p$ 上的延展 $p$ 维算符 $\cO$ 上。}
\end{figure}

\para{通过环绕的作用。} 描述作用 (2.58) 的一种等价方式是
\begin{equation}
  U_g\left(S^{d-p-1}\right)\cO(M_p) = \phi_g\times\cO(M_p)\,, \tag{2.63}
\end{equation}
其中左边我们沿一个环绕 $M_p$ 中某点 $x$ 的小 $(d-p-1)$ 维球面 $S^{d-p-1}$ 插入算符 $U$，如图 6 所示。

% [Tutor Note: p.19 (2.63) 把 (2.58) 中的 $\phi(g)$ 写成了 $\phi_g$（下标形式），而 p.18 图 6 的图注又用回 $\phi(g)$；原文此处记号不统一，但两处指同一个数。译文照原文分别记为 $\phi_g$ 与 $\phi(g)$，未作统一，读者当作同一对象处理即可。]

\para{荷与 Pontryagin 对偶群。} 阿贝尔群 $G^{(p)}$ 的不可约表示等价于同态
\begin{equation}
  \phi:\,G^{(p)}\to U(1) \tag{2.64}
\end{equation}
这类同态常被称为\textbf{特征标}（characters），它们自身构成一个阿贝尔群，其乘法运算为
\begin{equation}
  (\phi\phi')(g) = \phi(g)\times\phi'(g)\in U(1)\,. \tag{2.65}
\end{equation}
所得的群记作
\begin{equation}
  \widehat{G}^{(p)} = \{\text{由同态 } G^{(p)}\to U(1) \text{ 构成的群}\}\,, \tag{2.66}
\end{equation}
并称为 $G^{(p)}$ 的 \textbf{Pontryagin 对偶群}。

用这一语言来说，一个（不可约）$p$ 维算符在 $p$-形式对称群 $G^{(p)}$ 下所携带的荷是 Pontryagin 对偶群 $\widehat{G}^{(p)}$ 的一个元素。让我们讨论几个例子：

\begin{itemize}[leftmargin=1.4em,itemsep=1pt,topsep=1pt]
\item 对 $G^{(p)}=\U(1)$，我们有 $\widehat{G}^{(p)}=\Z$，即由整数在加法下构成的群。如果我们把 $G^{(p)}$ 的元素写成
\begin{equation}
  g = \ee^{\ii\alpha}\,,\qquad \alpha\in[0,2\pi) \tag{2.67}
\end{equation}
% [原文 p.20]
那么可能的同态为
\begin{equation}
  \phi(g) = g^q = \ee^{\ii q\alpha}\in U(1)\,,\qquad q\in\Z\,. \tag{2.68}
\end{equation}
\item 对 $G^{(p)}=\Z_N$，我们有 $\widehat{G}^{(p)}=\Z_N$，即由整数模 $N$ 在加法下构成的群。如果我们把 $G^{(p)}$ 的元素写成
\begin{equation}
  g = \ee^{\frac{2\pi\ii\alpha}{N}}\,,\qquad \alpha\in\{0,1,\cdots N-1\}\,, \tag{2.69}
\end{equation}
那么可能的同态为
\begin{equation}
  \phi(g) = g^q = \ee^{\frac{2\pi\ii q\alpha}{N}}\,,\qquad q\in\{0,1,\cdots N-1\}\,. \tag{2.70}
\end{equation}
\item 推广上一例子，若 $G^{(p)}$ 是\emph{有限}阿贝尔群，则我们有
\begin{equation}
  \widehat{G}^{(p)}\cong G^{(p)} \tag{2.71}
\end{equation}
这实际上可以作为上一例子的推论导出。对有限阿贝尔群，我们有
\begin{equation}
  G^{(p)}\cong\prod_{i=1}^{n}\Z_{N_i}\,,\qquad N_i\in\N \tag{2.72}
\end{equation}
从而
\begin{equation}
  \widehat{G}^{(p)}\cong\prod_{i=1}^{n}\widehat{\Z}_{N_i}\cong\prod_{i=1}^{n}\Z_{N_i} \tag{2.73}
\end{equation}
\end{itemize}

% [Tutor Note: p.20 (2.71)--(2.73) 只对\emph{有限}阿贝尔群成立：对紧阿贝尔群 $\U(1)$ 它不成立——(2.68) 已给出 $\widehat{\U(1)}=\Z\ne \U(1)$。另一处可核对的是，$\widehat{\Z_N}\cong\Z_N$ 的这个同构并不典范（依赖把 $\eta:\Z_N\times\Z_N\to\R/\Z$ 取成哪一个双同态），所以 (2.71) 只写 $\cong$ 不写典范同构，而 (2.74) 才强调“典范”。]

\begin{notebox}{双重 Pontryagin 对偶}
Pontryagin 对偶的一个重要的、我们将于后文用到的性质是：取两次 Pontryagin 对偶等价于根本不取 Pontryagin 对偶。更确切地说，对任意群 $G^{(p)}$ 都存在一个典范同构
\begin{equation}
  \widehat{\widehat{G}}^{(p)}\cong G^{(p)} \tag{2.74}
\end{equation}
事实上，一个元素 $g\in G^{(p)}$ 定义了一个同态
\begin{equation}
  h_g:\,\widehat{G}^{(p)}\to U(1) \tag{2.75}
\end{equation}
其形式为
\begin{equation}
  h_g(\phi) = \phi(g)\in U(1)\,,\qquad \phi\in\widehat{G}^{(p)}\,. \tag{2.76}
\end{equation}
\end{notebox}

% [原文 p.21]
\para{特例：$\U(1)$ $p$-形式对称性。} 推广例 2.1，$\U(1)$ $p$-形式对称性的 Noether 流的连续性方程，在存在一个荷为 $q\in\Z$ 的 $p$ 维算符 $\cO(M_p)$ 时，被修改为
\begin{equation}
  \cO(M_p)\dd j^{d-p-1} = q\delta^{d-p}(M_p)\cO(M_p)\,, \tag{2.77}
\end{equation}
其中 $\delta^{d-p}(M_p)$ 是与 $M_p$ 上 delta 函数相联系的 $(d-p)$-形式。

利用这一事实、并遵循与推导 (2.37) 类似的步骤，我们可以算出 $p$-形式对称性与 $\cO(M_p)$ 的环绕作用 (2.63) 为
\begin{equation}
  U_g\left(S^{d-p-1}\right)\cO(M_p) = \exp\left(\ii q\alpha\right)\cO(M_p)\,. \tag{2.78}
\end{equation}

\begin{examplebox}{例 2.3：Maxwell 理论中的高阶形式荷}
现在让我们回到例 2.2 中开始研究的纯 $\U(1)$ 规范理论。正如那里所讨论的，该理论有两个 $\U(1)$ 高阶形式对称性，即电 1-形式对称性和磁 $(d-3)$-形式对称性。

在电 1-形式对称性下带荷的算符是 \textbf{Wilson 线算符} $\cW_q$，当它沿时空中的一条线 $L$ 插入时可以写成
\begin{equation}
  \cW_q(L) = \exp\left(2\pi\ii q\int_L A\right)\,. \tag{2.79}
\end{equation}
不同的 Wilson 线由参数 $q\in\Z$ 区分，该参数常被称作 Wilson 线的“电荷”（electric charge）。算符 $\cW_q(L)$ 可以理解为描述一个在 $\U(1)$ 规范群下电荷为 $q$ 的无限重探针粒子，其世界线为 $L$。

在存在这样一条 Wilson 线时，运动方程被修改为
\begin{equation}
  (\dd\star F)\cW_q(L) = q\delta^{d-1}(L)\cW_q(L)\,, \tag{2.80}
\end{equation}
与 (2.77) 比较可知，Wilson 线 $\cW_q(L)$ 在 $\U(1)$ 电 1-形式对称性下的荷为 $q$。换言之，Wilson 线的电荷与它的 1-形式荷相同。因此，$\star F$ 在一个与 Wilson 线环绕的 $(d-2)$ 维流形上的积分测量该 Wilson 线的电荷：这可以看作 Gauss 定律的一个陈述。

在磁 $(d-3)$-形式对称性下带荷的算符是 \textbf{'t Hooft 算符}，它们是时空中的余维-3 算符。这些算符也被称为\textbf{磁单极子算符}，特别是在 $d=3$ 时。不同的 't Hooft 算符 $\cH_q(M_{d-3})$ 由参数 $q\in\Z$ 区分，该参数常被称作 't Hooft 算符的“磁荷”。算符 $\cH_q(M_{d-3})$ 可以理解为描述一个在 $\U(1)$ 规范群下磁荷为 $q$ 的无限重探针磁单极子，其世界体为 $M_{d-3}$。

't Hooft 算符 $\cH_q(M_{d-3})$ 的定义是：从时空中挖去 $M_{d-3}$，并要求路径积分中的规范场在 $M_{d-3}$ 附近具有如下边界条件
% [原文 p.22]
\begin{equation}
  \int_{S^2} F = q\,, \tag{2.81}
\end{equation}
其中 $S^2$ 是时空中与 $M_{d-3}$ 环绕的一个小球面。Bianchi 恒等式在存在 $\cH_q(M_{d-3})$ 时被修改：
\begin{equation}
  (\dd F)\cH_q(M_{d-3}) = q\delta^3(M_{d-3})\cH_q(M_{d-3})\,, \tag{2.82}
\end{equation}
与 (2.77) 比较可知，'t Hooft 算符 $\cH_q(M_{d-3})$ 在 $\U(1)$ 磁 $(d-3)$-形式对称性下的荷为 $q$。换言之，'t Hooft 算符的“磁荷”与它的 $(d-3)$-形式荷相同。
\end{examplebox}

% [Tutor Note: 例 2.3 方框在原文跨 pp.21--22 断开，此处按 STYLE\_GUIDE 第 7 节合并为一个 examplebox，编号 2.3 与内容均未拆分。另可核对：(2.81) 中的 $S^2$ 始终是 2 维球面（与 $M_{d-3}$ 环绕），因此 (2.81) 只在 $d\ge3$ 时有意义；$d=3$ 时 $M_{d-3}$ 退化为点，正是磁单极子的通常情形。]

\subsection{离散（高阶形式）规范理论}

在上一小节中，我们讨论了 Maxwell 理论中出现的电与磁高阶形式对称性。这两种高阶形式对称性都是连续的。本小节我们讨论一类展现出离散高阶形式对称性的理论。

我们要研究的这类理论是离散高阶形式规范理论，它们是拓扑量子场论（TQFT）的例子。一个 $d$ 维时空中的离散高阶形式规范理论具有作用量
\begin{equation}
  2\pi\int a_p\cup_\eta\delta b_{d-p-1}\,, \tag{2.83}
\end{equation}
其中
\begin{itemize}[leftmargin=1.4em,itemsep=1pt,topsep=1pt]
\item $a_p$ 是 $p$-形式场在时空上的一个离散版本，
\item $b_{d-p-1}$ 是 $(d-p-1)$-形式场在时空上的一个离散版本，
\item $\delta$ 是外微分的一个离散版本，且
\item $\cup_\eta$ 是楔积的一个离散版本，它被称为\emph{杯积}（cup product）。
\end{itemize}
对 $p=1$，我们有通常规范联络 $a_1$ 的离散类比物，因此相应的离散高阶形式规范理论就直接称为\textbf{离散规范理论}。

\para{离散规范场的上链描述。} 为了具体描述这些离散规范场，人们使用相对于时空的一个三角剖分而定义的上链（cochains）语言。在这一语言下，上面的各个要点成为：

\begin{itemize}[leftmargin=1.4em,itemsep=1pt,topsep=1pt]
\item $a_p$ 是一个 $G^{(p)}$ 取值的 $p$-上链，这意味着 $a_p$ 给三角剖分中的每个 $p$-单形指定阿贝尔群 $G^{(p)}$ 的一个元素。注意，对 $p=0$，取 $G^{(p)}=G^{(0)}$ 为阿贝尔群是我们在本小节所作的一个假设。有限非阿贝尔群的离散规范理论一般不允许形如 (2.83) 的描述。
% [原文 p.23]（本片只收尾这一跨页句；p.23 的其余内容见 tr01c。下面闭合列表，tr01c 以同款 itemize 续排）
\end{itemize}
